\chapter{Review of Literature}
\label{sec:2}

\section{Introduction}
\label{sec:lit_intro}

Research in automated safety detection using wearable and smartphone sensors has grown rapidly over the past two decades. The literature spans three overlapping areas relevant to this project: human fall detection, vehicular crash detection, and edge deployment of neural networks. This chapter surveys key contributions in each area and identifies the gaps that this project addresses.

\section{Detailed Review}
\label{sec:lit_detailed}

\subsection{Human Fall Detection}

Early systems for fall detection relied on threshold-based rules applied to accelerometer magnitude. Kangas et al. \citep{kangas2008} demonstrated that peak acceleration alone could flag falls with sensitivity above 95\%, making threshold methods attractive for simple wearable devices. In practice, however, these detectors generate frequent false positives during vigorous daily activities such as sitting down sharply, jumping, or running. The core weakness of a fixed threshold is that it cannot adapt to variation in sensor placement, individual body composition, or ambient vibration.

Deep learning changed this picture substantially. Musci et al. \citep{musci2020} applied Long Short-Term Memory (LSTM) networks to wearable IMU data and reported near-perfect sensitivity on several benchmark datasets. The recurrent architecture captures long-range temporal dependencies that threshold methods miss. Nho et al. \citep{nho2021} extended this line with a bidirectional LSTM that reached 99.1\% accuracy on the SisFall dataset. Despite these impressive numbers, LSTMs process each timestep sequentially, which creates a latency bottleneck on low-power edge hardware and raises power consumption during continuous streaming.

Convolutional approaches offer a parallelisable alternative. Wang et al. \citep{wang2020cnn} showed that 1D-CNNs applied to 50 Hz accelerometer streams can capture the localised temporal patterns that characterise fall events, using far fewer sequential operations than LSTMs. The weight-sharing property of convolutional filters also reduces the parameter count significantly.

\subsection{Vehicular Crash Detection}

Most crash detection research has focused on dedicated hardware embedded in vehicles rather than on the smartphones passengers already carry. Bhatt et al. \citep{bhatt2017} proposed a rule-based detector using accelerometer magnitude thresholds and reported 92\% accuracy. Under real driving conditions this approach proves brittle: hard braking, potholes, and dropping the phone can all trigger the same alarm as a genuine crash.

Jiang et al. \citep{jiang2020} improved on threshold methods by applying machine learning classifiers to naturalistic driving data collected from instrumented test vehicles. Their system reduced false alarms considerably, but it depends on specialised vehicle-embedded sensors that ordinary passengers do not have access to. More recent work has applied deep learning to smartphone IMUs for crash detection \citep{smartphone_crash2024}, but even these approaches require real labeled crash data for training. Collecting such data at scale without exposing participants to actual crashes is practically infeasible.

\subsection{Edge Deployment of Neural Networks}

Compressing neural networks for constrained hardware has matured rapidly. TensorFlow Lite supports weight quantization, pruning, and operator fusion that together make deep models viable on microcontrollers \citep{tflite_micro}. Surveys on TinyML confirm that safety-critical inference on IoT devices is no longer a research curiosity but a practical reality \citep{tinyml2023, edge_safety2024}.

David et al. \citep{david2021} demonstrated that ARM Cortex-M processors can run neural networks occupying as few as 20 KB of flash. Processing sensor data locally eliminates cloud round-trip latency, works in areas with no network coverage, and never exposes raw motion streams to a remote server.

\section{Gap Identified}
\label{sec:gap}

A review of the literature reveals the following gaps that this project addresses:

\begin{enumerate}
    \item \textbf{Dual-mode detection on a single sensor}: No existing smartphone-based system simultaneously detects both human in-vehicle falls and vehicular crashes using a single 6-axis IMU pipeline and a unified inference engine. Prior works address these as separate problems.
    \item \textbf{Synthetic crash data}: Real labeled crash recordings are scarce and largely proprietary. No prior work has demonstrated training a smartphone crash detector entirely on physics-derived synthetic data and validating it against physical sensor behaviour.
    \item \textbf{Integrated SOS with GPS}: Existing systems detect events but stop short of triggering an automated, GPS-tagged emergency alert. This project integrates detection and alerting end-to-end.
    \item \textbf{Edge-optimised streaming pipeline}: Most prior work evaluates models offline. This project validates the full 50 Hz streaming pipeline on consumer hardware with measured sub-3 ms inference latency.
\end{enumerate}
